\subsection{BÀI TẬP TRẮC NGHIỆM}
\ind{PHẦN I.} \inden{Câu trắc nghiệm nhiều phương án lựa chọn. Mỗi câu hỏi học sinh chỉ chọn một phương án.}\\
\setcounter{ex}{0}
\Opensolutionfile{ans}[ans/1D7-B1-1]
\begin{ex}%[Dự án chuyển word 10-11, Chu Đức Minh]%[1D5B1-1]
	Tính số gia của hàm số $y=x^2+2$ tại điểm $x_0=2$ ứng với số gia $\Delta x=1$.
	\choice
	{$\Delta y=13$}
	{$\Delta y=9$}
	{\True  $\Delta y=5$}
	{$\Delta y=2$}
	\loigiai{
		Ta có $\Delta y=f\left( x_0+\Delta x \right)-f\left( x_0 \right)=f\left( 2+1 \right)-f\left( 2 \right)=f\left( 3 \right)-f\left( 2 \right)$
		$=\left( 3^2+2 \right)-\left( 2^2+2 \right)=5.$
	}
\end{ex}

\begin{ex}%[Dự án chuyển word 10-11, Chu Đức Minh]%[1D5B1-1]
	Tính số gia của hàm số $y=x^3+x^2+1$ tại điểm $x_0$ ứng với số gia $\Delta x=1$.
	\choice
	{$\Delta y=3x_0^2+5x_0+3$}
	{$\Delta y=2x_0^3+3x_0^2+5x_0+2$}
	{\True $\Delta y=3x_0^2+5x_0+2$}
	{$\Delta y=3x_0^2-5x_0+2$}
	\loigiai{
		Ta có 
		\begin{eqnarray*}
			\Delta y &=&f\left( x_0+\Delta x \right)-f\left( x_0 \right)=f\left( x_0+1 \right)-f\left( x_0 \right)\\
			&=&\left[\left( x_0+1 \right)^3+\left(x_0+1\right)^2+1 \right]-\left[ x_0^3+x_0^2+1 \right]\\
			&=&3x_0^2+5x_0+2.
		\end{eqnarray*}
	}
\end{ex}

\begin{ex}%[Dự án chuyển word 10-11, Chu Đức Minh]%[1D5B1-1]
	Tính số gia của hàm số $y=\dfrac{x^2}{2}$ tại điểm $x_0=-1$ ứng với số gia $\Delta x$. 
	\choice
	{\True  $\Delta y=\dfrac{1}{2}\left( \Delta x \right)^2-\Delta x$}
	{$\Delta y=\dfrac{1}{2}\left[ \left( \Delta x \right)^2-\Delta x \right]$}
	{$\Delta y=\dfrac{1}{2}\left[\left( \Delta x \right)^2+\Delta x \right]$}
	{$\Delta y=\dfrac{1}{2}\left( \Delta x \right)^2+\Delta x$} 
	\loigiai{
		Ta có 
		\begin{eqnarray*}
			\Delta y&=& f\left( x_0+\Delta x \right)-f\left( x_0 \right)=f\left( -1+\Delta x \right)-f\left( -1 \right)\\
			&=& \dfrac{\left( -1+\Delta x \right)^2}{2}-\dfrac{1}{2}=\dfrac{1-2\Delta x+\left( \Delta x \right)^2}{2}-\dfrac{1}{2}\\
			&=&\dfrac{1}{2}\left( \Delta x \right)^2-\Delta x.
		\end{eqnarray*}
	}
\end{ex}

\begin{ex}%[Dự án chuyển word 10-11, Chu Đức Minh]%[1D5B1-1]
	Tính số gia của hàm số $y=x^2-4x+1$ tại điểm $x_0$ ứng với số gia $\Delta x$. 
	\choice
	{\True  $\Delta y=\Delta x\left( \Delta x+2x_0-4 \right)$}
	{$\Delta y=2x_0+\Delta x$}
	{$\Delta y=\Delta x\left( 2x_0-4\Delta x \right)$}
	{$\Delta y=2x_0-4\Delta x$} 
	\loigiai{
		Ta có 
		\begin{eqnarray*}
			\Delta y&=&f\left( x_0+\Delta x \right)-f\left( x_0 \right)\\
			&=&\left[\left( x_0+\Delta x \right)^2-4\left( x_0+\Delta x \right)+1 \right]-\left[ x_0^2-4x_0+1 \right]\\
			&=&\Delta x\left( \Delta x+2x_0-4 \right).
		\end{eqnarray*}
	}
\end{ex}

\begin{ex}%[Võ Đông Phước]%[1D5B1]
	Trong các hàm số dưới đây, hàm số nào có $f'(-3)=7$?
	\choice{$f(x)=\dfrac{1-2x}{x+4}$}
	{$f(x)=\sqrt{x^2+40}$}
	{$f(x)=x^2-x+8$}
	{\True $f(x)=\dfrac{3x-1}{x+2}$}
\end{ex}

\begin{ex}%[Bùi Sang Thọ]%[1D5B1]
	Đạo hàm của hàm số $y=\dfrac{1}{2x-1}$ tại điểm $x_0=1$ là
	\choice
	{$2$}
	{\True $-2$}
	{$-1$}
	{$1$}
\end{ex}

\begin{ex}%[Võ Đông Phước]%[1D5B1]
	Cho hàm số $f(x)=(m^3-3m^2+m)x+2017$. Tìm m để $f'(2)=3$.
	\choice{ $m=-1$}
	{ $m=1$}
	{\True $m=3$}
	{ $m=0$}
\end{ex}

\begin{ex}%[Hoàng Trình]%[1D5B1]
	Đạo hàm của hàm số $f(x)=3x-1$ tại $x_0=1$ là
	\choice
	{$0$}
	{$2$}
	{$1$}
	{\True $3$}
\end{ex}

\begin{ex}%[Hoàng Trình]%[1D5B1]
	Hệ số góc của tiếp tuyến với đồ thị hàm số $f(x)=-x^3$ tại điểm $M(-2; 8)$ là
	\choice
	{$12$}
	{\True $-12$}
	{$192$}
	{$-192$}
\end{ex}

\begin{ex}%[Võ Đông Phước]%[1D5B1]
	Tìm phương trình tiếp tuyến của đồ thị $(C)$: $y=\dfrac{1}{x}$ tại điểm có hoành độ bằng $-1$.
	\choice{\True $y=-x-2$}
	{$y=x-2$}
	{$y=-x+1$}
	{$y=-x-1$}
\end{ex}

\begin{ex}%[Dự án chuyển word 10-11, Chu Đức Minh]%[1D5B2-2]
	Viết phương trình tiếp tuyến của đường cong $y=x^3$ tại điểm $\left( -1;-1 \right).$
	\choice
	{$y=-3x-4$ }
	{$y=-1$}
	{$y=3x-2$}
	{\True $y=3x+2$}
	\loigiai{
		\begin{itemize}
			\item Ta tính được $k=y'(-1)=3$. 
			\item Ta có  $\heva{& x_0=-1 \\ & y_0=-1 \\ & k=3.}$
			\item Suy ra phương trình tiếp tuyến $y+1=3(x+1)\Leftrightarrow y=3x+2$. 
		\end{itemize}		
	}
\end{ex}

\begin{ex}%[Dự án chuyển word 10-11, Chu Đức Minh]%[1D5B2-2]
	Viết phương trình tiếp tuyến của đường cong $y=\dfrac{1}{x}$ tại điểm có hoành độ bằng $-1$.
	\choice
	{\True $x+y+2=0$}
	{$y=x+2$ }
	{$y=x-2$}
	{$y=-x+2$}
	\loigiai{
		\begin{itemize}
			\item Ta tính được $k=y'(-1)=-1$. 
			\item Với $x_0=-1\Rightarrow y_0=-1$.
			\item Ta có $\heva{& x_0=-1 \\ 
				& {{y}_0}=-1 \\ 
				& k=-1 \\ 
			}$. 
			Suy ra phương trình tiếp tuyến $y+1=-1(x+1)\Leftrightarrow y=-x-2$.
		\end{itemize}
	}
\end{ex}

\begin{ex}%[1D5B2-2]
	Viết phương trình tiếp tuyến của đường cong $y=x^3$ tại điểm có tung độ bằng $8$.
	\choice
	{$y=8$}
	{$y=-12x+16$}
	{$y=12x-24$}
	{\True $y=12x-16$}
	\loigiai{
		\begin{itemize}
			\item Với $y_0=8\Rightarrow x_0=2$.
			\item Ta tính được $k=y'\left( 2 \right)=12$. 
			\item Ta có $\heva{& x_0=2 \\& {{y}_0}=8 \\& k=12 \\ } \Rightarrow$ phương trình tiếp tuyến $y-8=12\left( x-2 \right)\Leftrightarrow y=12x-16.$
		\end{itemize}
	}
\end{ex}

\begin{ex}%[1D5B2-3]
	Viết phương trình tiếp tuyến của đường cong $y=\dfrac{1}{x}$ biết hệ số góc của tiếp tuyến bằng $-\dfrac{1}{4}$. 
	\choice
	{$x+4y-1=0, x+4y+1=0$}
	{\True $x+4y-4=0, x+4y+4=0$}
	{$y=-\dfrac{1}{4}x-4, y=-\dfrac{1}{4}x+4$}
	{$y=-\dfrac{1}{4}x$}
	\loigiai{
		\begin{itemize}
			\item Gọi $M\left(x_0;y_0\right)$ là tọa độ tiếp điểm. 
			Ta tính được $k=y'\left(x_0 \right)=-\dfrac{1}{x_0^2}$. 
			\item Theo giả thiết ta có $k=-\dfrac{1}{4}\Leftrightarrow -\dfrac{1}{x_0^2}=-\dfrac{1}{4}\Leftrightarrow x_0^2=4\Leftrightarrow x_0=\pm 2.$
			\item Với $x_0=2$, ta có $y_0=\dfrac{1}{2}$
			nên phương trình tiếp tuyến cần tìm là $y=-\dfrac{1}{4}\left(x-2 \right)+\dfrac{1}{2}\Leftrightarrow x+4y-4=0$. 
			\item Với $x_0=-2$, ta có $y_0=-\dfrac{1}{2}$ nên phương trình tiếp tuyến cần tìm là $$y=-\dfrac{1}{4}\left( x+2 \right)-\dfrac{1}{2}\Leftrightarrow x+4y+4=0.$$
		\end{itemize}
	}
\end{ex}

\begin{ex}%[Dự án chuyển word 10-11, Chu Đức Minh]%[1D5B2-6]
	Một chất điểm chuyển động theo phương trình $s(t)=t^2$, trong đó $t>0,$ $t$ tính bằng giây và $s(t)$ tính bằng mét. Tính vận tốc của chất điểm tại thời điểm $t=2$ giây.
	\choice
	{$2$ m/s }
	{$3$ m/s }
	{\True $4$ m/s }
	{$5$ m/s}
	\loigiai{
		Ta tính được $s'(t)=2t.$\\
		Vận tốc của chất điểm $v(t)=s'(t)=2t\Rightarrow v\left(2\right)=2\cdot 2=4$ m/s. } 
\end{ex}

\begin{ex}%[Bùi Sang Thọ]%[1D5B1]
	Một vật rơi tự do theo phương trình $s=\dfrac{1}{2}gt^2$, trong đó $g\approx 9,8\;\text{m/s}^2$ là gia tốc trọng trường. Tìm vận tốc tức thời của chuyển động tại thời điểm $t=5 $s.
	\choice
	{\True $49$ m/s}
	{$122,5$ m/s}
	{$50$ m/s}
	{$123$ m/s}
\end{ex}

\Closesolutionfile{ans}

\ind{PHẦN II.} \inden{Câu trắc nghiệm đúng sai. Trong mỗi ý a), b), c), d) ở mỗi câu, học sinh chọn đúng hoặc sai.}\\
\Opensolutionfile{ans}[ans/1D7-B1-2]

\begin{ex}%[1D7V1-1]
	Dùng định nghĩa để tính đạo hàm của hàm số $f(x)=2x^3$. Xét tính đúng sai của các mệnh đề sau.
	\choiceTF 
	{\True Với bất kì $x_0$ thì $f'(x_0)=\lim\limits_{x\to x_0}\dfrac{f(x)-f(x_0)}{x-x_0}$}
	{$f'(1)=-6$}
	{\True $f'(0)=0$}
	{\True $f'(2)=24$}
	\loigiai{
		\begin{itemchoice}
			\itemch Đúng. Hàm số $f(x)=2x^3$ xác định trên $\mathbb{R}$ và \begin{align*}
				\lim\limits_{x\to x_0}\dfrac{f(x)-f(x_0)}{x-x_0}=\lim\limits_{x\to x_0}\dfrac{2x^3-2x_0^3}{x-x_0}=\lim\limits_{x\to x_0}\left(2x^2+2x_0x+2x_0^2\right)=6x_0^2.
			\end{align*} Vậy $f'(x_0)=\lim\limits_{x\to x_0}\dfrac{f(x)-f(x_0)}{x-x_0}=6x_0^2$.
			\itemch Sai. Vì $f'(1)=6\cdot 1^2=6\neq -6$.
			\itemch Đúng. Vì $f'(0)=6\cdot 0^2=0$.
			\itemch Đúng. Vì $f'(2)=6\cdot 2^2=24$.
		\end{itemchoice}
	}
\end{ex}

\begin{ex}%[1D7H1-1]
	Dùng định nghĩa để tính đạo hàm của hàm số $f(x)=\dfrac{2}{1-x}$ với $x \neq 1$. Xét tính đúng sai của các mệnh đề sau.
	\choiceTF 
	{\True Với bất kì $x_0 \neq 1$, ta có $f'(x_0)=\lim\limits_{x\to x_0}\dfrac{2}{(1-x)\left(1-x_0\right)}$}
	{\True $f'(2)=2$}
	{$f'(3)=\dfrac{1}{3}$}
	{\True $f'(2)+f'(3)=\dfrac{5}{2}$}
	\loigiai{
		\begin{itemchoice}
			\itemch Đúng. Hàm số $f(x)=\dfrac{2}{1-x}$ có tập xác định $\mathscr{D}=\mathbb{R}\setminus\{1\}$. Với $x\neq 1$ thì \begin{align*}
				\lim\limits_{x\to x_0}\dfrac{f(x)-f(x_0)}{x-x_0}&=\lim\limits_{x\to x_0}\dfrac{\dfrac{2}{1-x}-\dfrac{2}{1-x_0}}{x-x_0}\\
				&=\lim\limits_{x\to x_0}\dfrac{2(x-x_0)}{(x-x_0)(1-x)(1-x_0)}\\
				&=\lim\limits_{x\to x_0}\dfrac{2}{(1-x)(1-x_0)}.
			\end{align*}
			Vậy $f'(x_0)=\lim\limits_{x\to x_0}\dfrac{2}{(1-x)(1-x_0)}$.
			\itemch Đúng. Ta có $f'(x_0)=\lim\limits_{x\to x_0}\dfrac{2}{(1-x)(1-x_0)}=\dfrac{2}{(1-x_0)^2}$ với $x_0\neq 1$. Do đó $f'(2)=2$.
			\itemch Sai. Vì $f'(3)=\dfrac{2}{(1-3)^2}=\dfrac{1}{2}\neq \dfrac{1}{3}$.
			\itemch Đúng. Ta có $f'(2)+f'(3)=2+\dfrac{1}{2}=\dfrac{5}{2}$.
		\end{itemchoice}
	}
\end{ex}
\Closesolutionfile{ans}

\ind{PHẦN III.} \inden{Câu trắc nghiệm trả lời ngắn.}\\
\Opensolutionfile{ans}[ans/1D7-B1-3]
\begin{ex}%[Dự án chuyển word 10-11, Chu Đức Minh]%[1D5B2-6]
	Một viên đạn được bắn lên cao theo phương trình $s(t)=196t-4{,}9t^2$ trong đó $t>0,$ $t$ tính bằng giây kể từ thời điểm viên đạn được bắn lên cao và $s(t)$ là khoảng cách của viên đạn so với mặt đất được tính bằng mét. Tại thời điểm vận tốc của viên đạn bằng $0$ thì viên đạn cách mặt đất bao nhiêu mét?\\
	\shortans[oly]{$1960$}
	\loigiai{
		\begin{itemize}
			\item Vận tốc của viên đạn $v(t)=s'(t)=196-9{,}8t$. 
			\item Ta có $v(t)=0\Leftrightarrow 196-9{,}8t=0\Leftrightarrow t=20$.
			\item Khi đó viên đạn cách mặt đất một khoảng $h=s\left(20 \right)=196\cdot 20-4{,}9 \cdot 20^2=1960$ m. 
		\end{itemize}
	}
\end{ex}


\begin{ex}%[Dự án chuyển word 10-11, Chu Đức Minh]%[1D5B2-6]
	Vận tốc của một chất điểm chuyển động được biểu thị bởi công thức $v(t)=8t+3t^2$, trong đó $t>0,$ $t$ tính bằng giây và $v(t)$ tính bằng mét/giây. Tìm gia tốc (m/s$^2$) của chất điểm tại thời điểm mà vận tốc chuyển động là $11$ mét/giây.\\
	\shortans[oly]{$14$}
	\loigiai{
		\begin{itemize}
			\item Ta có $a(t) = v'(t)=8+6t$. 
			\item Ta có $v(t)=11\Leftrightarrow 8t+3t^2=11\Leftrightarrow t=1 (t>0)$. 
			\item Gia tốc của chất điểm $a(t)=v'(t)=8+6t\Rightarrow a(1)=v'(1)= 8+6 \cdot 1=14$ m/s$^2$. 
		\end{itemize}	
	}
\end{ex}

\begin{ex}%[Dự án chuyển word 10-11, Chu Đức Minh]%[1D5B2-6]
	Một vật rơi tự do theo phương trình $s=\dfrac{1}{2}gt^2$, trong đó $g=9{,}8$  m/s $^2$ là gia tốc trọng trường. Tìm vận tốc trung bình (m/s) của chuyển động trong khoảng thời gian từ  $t$ ($t=5$s) đến $t+\Delta t$ với $\Delta t=0{,}001$s (kết quả làm tròn đến hàng đơn vị).
	\\
	\shortans[oly]{$49$}
	\loigiai{
		Ta có $v_{\text{tb}} = \dfrac{s(t + \Delta t) - s(t)}{\Delta t} = \dfrac{\dfrac{1}{2}g(t + \Delta t)^2 - \dfrac{1}{2}gt^2}{\Delta t} = gt + \dfrac{1}{2}g\Delta t = 49{,}0049$ m/s. 
	}
\end{ex}

\begin{ex}%[1D7H1-4]
	Cho biết điện lượng truyền trong dây dẫn theo thời gian biểu thị bởi hàm số $Q(t)=2t^2+t$, trong đó $t$ được tính bằng giây và $Q$ được tính theo Cu-lông. Tính cường độ dòng điện tại thời điểm $t=4$ (s).\\
	\shortans[oly]{$17$}
	\loigiai{Tại thời điểm $t_0$ ta có \begin{align*}
			\lim\limits_{t\to t_0}\dfrac{Q(t)-Q(t_0)}{t-t_0}&=\lim\limits_{t\to t_0}\dfrac{2t^2+t-\left(2t_0^2+t_0\right)}{t-t_0}\\
			&=\lim\limits_{t-t_0}\dfrac{(t-t_0)\left[2(t+t_0)+1\right]}{t-t_0}\\
			&=\lim\limits_{t\to t_0}2(t+t_0)+1=4t_0+1.
		\end{align*} 
		Do đó, $Q'(t_0)=4t_0+1$.\\
		Vậy
		cường độ dòng điện tại thời điểm $t=4$ (s) là $Q'(4)=4\cdot 4+1=17$ (C/s).
	}
\end{ex}

\begin{ex}%[1D7V1-4]
	Một vật được phóng theo phương thẳng đứng lên trên từ mặt đất, biết độ cao $h$ của nó (tính bằng mét) sau $t$ giây được cho bởi phương trình $h(t)=24{,}5t-4{,}9t^2$. Tìm vận tốc của vật khi nó chạm đất.\\
	\shortans[oly]{$24{,}5$}
	\loigiai{Ta có \begin{align*}
			\lim\limits_{t\to t_0}\dfrac{h(t)-h(t_0)}{t-t_0}&=\lim\limits_{t\to t_0}\dfrac{24{,5}t-4{,}9t^2-\left(24{,}5t_0-4{,}9t_0^2\right)}{t-t_0}\\
			&=\lim\limits_{t\to t_0}\dfrac{(t-t_0)\left[24{,}5-4{,}9(t+t_0)\right]}{t-t_0}\\
			&=\lim\limits_{t\to t_0}\left[24{,}5-4{,}9(t+t_0)\right]\\
			&=24{,}5-9{,}8t_0.
		\end{align*}
		Do đó vận tốc của vật tại thời điểm $t_0$ là $v(t_0)=h'(t_0)=24{,}5-9{,}8t_0$.\\
		Khi vật chạm đất thì $h=0 \Leftrightarrow 24{,}5t-4{,}9t^2=0 \Rightarrow t=5$.\\
		Vậy tốc độ của vật tại thời điểm nó chạm đất $t=5$ là $$\left|v(5)\right|=\left|h'(5)\right|=|24{,}5-9{,}8\cdot 5|=24{,}5\  (\mathrm{m}/\mathrm{s}).$$
	}
\end{ex}

\begin{ex}%[1D7V1-4]
	Một quả bóng được thả rơi tự do từ đài quan sát trên sân thượng của toà nhà Landmark 81 (Thành phố Hồ Chí Minh) cao $461{,}34$ m xuống mặt đất, với phương trình chuyển động $s(t)=4{,}9t^2$. Tính vận tốc của quả bóng khi nó chạm đất, bỏ qua sức cản không khí. (Đơn vị m/s, kết quả gần đúng làm tròn đến hàng phần chục).\\
	\shortans[oly]{$ 95{,}1$}
	\loigiai{
		Với bất kì $t_0$, ta có
		$\lim\limits_{t\to t_0} \dfrac{s(t)-s\left(t_0\right)}{t-t_0}=\lim\limits_{t\to t_0} \dfrac{4{,}9 t^2-4{,}9 t_0^2}{t-t_0}=\lim\limits_{t\to t_0} 4{,}9\left(t+t_0\right)=9{,}8 t_0$.\\
		Do đó, vận tốc của quả bóng tại thời điểm $t$ là $v(t)=s'(t)=9{,}8t$ m/s.\\
		Mặt khác, vì chiều cao của toà tháp là $461{,}3$ m nên quả bóng sẽ chạm đất tại thời điểm $t_1$. Từ đó, ta có $4{,}9 t_1^2=461{,}3 \Leftrightarrow t_1=\sqrt{\dfrac{461{,}3}{4{,}9}}$ (giây).\\
		Vậy vận tốc của quả bóng khi nó chạm đất là
		$v\left(t_1\right)=9{,}8 t_1=9{,}8 \cdot \sqrt{\dfrac{461{,}3}{4{,}9}} \approx 95{,}1$ m/s.
	}
\end{ex}

\begin{ex}%[1D7C1-5]
	\immini[thm]{Người ta xây dựng một cây cầu vượt giao thông hình parabol nối hai điểm có khoảng cách là $400$ m. Độ dốc của mặt cầu không vượt quá $10^\circ$ (độ dốc tại một điểm được xác định bởi góc giữa phương tiếp xúc với mặt cầu và phương ngang). Tính chiều cao giới hạn từ đỉnh cầu đến mặt đường (làm tròn kết quả đến chữ số thập phân thứ nhất).
}{
\begin{tikzpicture}[scale=1.5, font=\footnotesize, line join=round, line cap=round,>=stealth]
	\draw[smooth,samples=100,domain=-2:2] plot(\x,{-0.25*(\x)^2});
	\draw[smooth,samples=100,domain=-1.7:-0.3] plot(\x,{-0.5*(-1.2)*(\x+1.2)-0.36});
	\draw (-2,-1)--(2,-1)node[below,pos=0.5]{$400$ m};
	\path (-1.2,-0.36) coordinate (P) (-0.6,0) coordinate (a) (0,-0.36) coordinate (J);
	\draw[dashed] (P)--(-0.5,-0.36) (0,-1)--(0,0)node[right,pos=0.5]{$?$};
	\foreach \x/\y/\l/\g  in {-2/-1/A/-90,2/-1/B/-90,-1.2/-0.36/P/150}{\fill[black] (\x,\y) circle(1pt)+(\g:0.2) node{$\l$};}
	\draw pic[draw, angle eccentricity=1.2, angle radius=0.5cm]{angle=J--P--a};
\end{tikzpicture}}

	\shortans[oly]{$17{,}6$}
	\loigiai{
		\begin{center}
			\begin{tikzpicture}[scale=1.5, font=\footnotesize, line join=round, line cap=round,>=stealth]
				\def \xmin{-2.5};
				\def \xmax{2.7};
				\def \ymin{-1.5};
				\def \ymax{1.0};
				\draw[->] (\xmin, 0.) -- (\xmax,0.) node[anchor=north] {$x$};
				\draw[->] (0.,\ymin) -- (0.,\ymax) node[anchor=west] {$y$};
				\clip(\xmin+0.1,\ymin+0.1) rectangle (\xmax-0.1,\ymax-0.1);
				\draw[smooth,samples=100,domain=-2:2] plot(\x,{-0.25*(\x)^2});
				\draw[smooth,samples=100,domain=-1.7:-0.3] plot(\x,{-0.5*(-1.2)*(\x+1.2)-0.36});
				\draw[dashed] (-2,0) node[above]{$-200$} |- (0,-1)node[]{}-|(2,0)node[above]{$200$} (-1.2,-0.36)--(-0.5,-0.36);
				\path (-1.2,-0.36) coordinate (P) (-0.6,0) coordinate (a) (0,-0.36) coordinate (J);
				\foreach \x/\y/\l/\g  in {0/0/O/45,-2/-1/A/-90,2/-1/B/-90,-1.2/-0.36/P/150,0/-1/I/-45}{\fill[black] (\x,\y) circle(1pt)+(\g:0.2) node{$\l$};}
				\draw pic[draw, angle eccentricity=1.2, angle radius=0.5cm]{angle=J--P--a};
			\end{tikzpicture}
		\end{center}
		Chọn hệ trục toạ độ như hình vẽ, sao cho đỉnh cầu là gốc tọa độ và mặt cắt của cây cầu có hình dạng parabol $y=-a x^2$ (với $a$ là hằng số dương).\\
		Hệ số góc của tiếp tuyến của parabol bằng $k=y'(x_0)=-2 a x_0,-200 \leq x_0 \leq 200$.\\
		Hệ số góc xác định độ dốc của mặt cầu (độ dốc dương) là $|k|=2 a|x| \leq 400 a$.\\
		Vì độ dốc của mặt cầu không vượt quá $10^\circ$ nên ta có
		$400 a \leq \tan 10^\circ \Leftrightarrow a \leq \dfrac{\tan 10^\circ}{400}$.\\
		Chiều cao giới hạn từ đỉnh cầu đến mặt đường là đoạn $O I$, cũng chính là độ lớn của tung độ điểm $B$ khi $a$ đạt giá trị lớn nhất.\\
		Do đó, $O I=\left|-a \cdot 200^2\right|\approx 17{,}6$ (m).\\
		Vậy chiều cao giới hạn từ đỉnh cầu đến mặt đường là $17{,}6$ m.
	}
\end{ex}

\Closesolutionfile{ans}

